\chapter{Separación de cerrados y extensión de Tietze}\label{chap:12}
\section{Extender valores sin crear discontinuidades}
Una función definida en una parte de un espacio puede contener información que deseamos conservar al ampliar el dominio. La continuidad uniforme permite extenderla a una clausura cuando el espacio de llegada es completo. El teorema de Tietze resuelve otro problema: una función real continua sobre un \emph{cerrado} de un espacio métrico se extiende a todo el espacio, sin exigir continuidad uniforme. La construcción combina funciones que separan cerrados y una serie de correcciones de tamaño decreciente.

\section{Separar cerrados mediante funciones}
\begin{teorema}[Lema de Urysohn para espacios métricos]\label{ch12:urysohn}
Si $A,B$ son cerrados disjuntos de un espacio métrico $X$, existe una función continua $u:X\to[0,1]$ tal que $u|_A=0$ y $u|_B=1$.
\end{teorema}
\begin{proof}
Si $A=\varnothing$, sirve $u=1$; si $B=\varnothing$ y $A\ne\varnothing$, sirve $u=0$. Cuando ambos son no vacíos, definimos
\[
 u(x)=\frac{\dist(x,A)}{\dist(x,A)+\dist(x,B)}.
\]
Las funciones de distancia son continuas, incluso $1$-Lipschitz. El denominador no se anula: si ambas distancias fueran cero, el carácter cerrado de los conjuntos obligaría a $x\in A\cap B$. Por tanto el cociente es continuo y está entre cero y uno. En $A$, el numerador vale cero; en $B$, la distancia a $B$ es cero y la de $A$ es positiva, por lo que el cociente vale uno. Esto también permite separar los cerrados por los abiertos $u^{-1}([0,1/3))$ y $u^{-1}((2/3,1])$, entendiendo los intervalos como abiertos relativos de $[0,1]$.
\end{proof}
\index{Urysohn, lema de}\index{separacion de cerrados@separación de cerrados}
Los cerrados pueden tener distancia mutua cero; no se exige una separación uniforme. Lo que se necesita es que, para cada punto individual, no se anulen simultáneamente las dos distancias.

\begin{lema}[Una corrección que reduce el error]\label{ch12:correccion}
Sea $V$ cerrado en $X$, y sea $g:V\to\R$ continua con $|g|\leq\alpha$, donde $\alpha>0$. Existe una función continua $h:X\to\R$ tal que
\[
 |h(x)|\leq\alpha/3\quad(x\in X),\qquad
 |g(x)-h(x)|\leq2\alpha/3\quad(x\in V).
\]
\end{lema}
\begin{proof}
Los conjuntos
\[
 A=\{x\in V:g(x)\leq-\alpha/3\},\qquad
 B=\{x\in V:g(x)\geq\alpha/3\}
\]
son cerrados en $V$ por continuidad, y por tanto cerrados en $X$ porque $V$ lo es. Son disjuntos. Sea $u$ la función del lema de Urysohn con $u=0$ en $A$ y $u=1$ en $B$, y pongamos $h=(\alpha/3)(2u-1)$. Entonces $|h|\leq\alpha/3$, $h=-\alpha/3$ en $A$ y $h=\alpha/3$ en $B$.

Si $x\in A$, los valores de $g(x)$ están entre $-\alpha$ y $-\alpha/3$, de modo que $|g(x)+\alpha/3|\leq2\alpha/3$. En $B$, $g(x)$ está entre $\alpha/3$ y $\alpha$, y $|g(x)-\alpha/3|\leq2\alpha/3$. En $V\setminus(A\cup B)$ se tiene $|g(x)|<\alpha/3$ y $|h(x)|\leq\alpha/3$, así que la desigualdad triangular da el mismo resultado. Estos tres casos cubren todo $V$.
\end{proof}

\section{Las herramientas sobre convergencia uniforme}
\begin{definicion}
Una sucesión $f_n:X\to\R$ converge \emph{uniformemente} a $f$ si
\[
 \forall\eps>0\ \exists N\in\N\ \forall n\geq N\ \forall x\in X,
 \qquad |f_n(x)-f(x)|<\eps.
\]
Una serie de funciones converge uniformemente si lo hace su sucesión de sumas parciales. El índice $N$ depende de $\eps$, pero no de $x$.
\end{definicion}
\index{convergencia!uniforme}\index{Weierstrass!criterio M}

\begin{teorema}[Criterio M de Weierstrass]\label{ch12:weierstrass}
Si $|f_n(x)|\leq M_n$ para todo $x\in X$, donde $M_n\geq0$ y $\sum_nM_n<\infty$, entonces $\sum_nf_n$ converge uniformemente en $X$.
\end{teorema}
\begin{proof}
Para cada punto $x$, la serie de valores absolutos está dominada por la serie numérica convergente, así que $f(x)=\sum_{n\geq1}f_n(x)$ existe en $\R$. Para todo $N$ y todo $x$,
\[
 \left|f(x)-\sum_{n=1}^Nf_n(x)\right|
 \leq\sum_{n>N}|f_n(x)|\leq\sum_{n>N}M_n.
\]
La última cola tiende a cero y no depende de $x$. Dado $\eps>0$, tomamos $N_0$ con $\sum_{n>N_0}M_n<\eps$; para todo $N\geq N_0$ y todo $x$, el error es menor que $\eps$. Ésta es la definición de convergencia uniforme.
\end{proof}

\begin{teorema}[Continuidad del límite uniforme]\label{ch12:limite-continuo}
Si $f_n:X\to\R$ son continuas y convergen uniformemente a $f$, entonces $f$ es continua.
\end{teorema}
\begin{proof}
Fijemos $x\in X$ y $\eps>0$. Elegimos un índice $N$ para el que $|f_N(z)-f(z)|<\eps/3$ para \emph{todo} $z\in X$. Usamos la continuidad de esta función fija $f_N$ en $x$ para elegir $\delta>0$ con $|f_N(y)-f_N(x)|<\eps/3$ cuando $d(x,y)<\delta$. Entonces
\[
 |f(y)-f(x)|\leq |f(y)-f_N(y)|+|f_N(y)-f_N(x)|+|f_N(x)-f(x)|<\eps.
\]
Así $f$ es continua en $x$. La elección uniforme del índice permite controlar a la vez el error en $x$ y en los puntos vecinos $y$.
\end{proof}

\section{Demostración completa del teorema de extensión}
\begin{teorema}[Tietze]\label{ch12:tietze}
Sea $V$ un cerrado de un espacio métrico $X$. Toda función continua $g:V\to\R$ tiene una extensión continua $f:X\to\R$. Si $g(V)\subset[a,b]$, puede elegirse una extensión con valores en el mismo intervalo.
\end{teorema}
\index{Tietze, teorema de}\index{extension de funciones@extensión de funciones}
\begin{proof}
Si $V=\varnothing$, una función constante resuelve el problema; para la afirmación del intervalo elegimos su valor en $[a,b]$. Si $X=V$, sirve $g$. Tratamos el caso restante.

\textbf{Caso acotado y simétrico.} Supongamos $|g|\leq\alpha$ para algún $\alpha>0$. Aplicamos el lema~\ref{ch12:correccion} a $g$ y obtenemos $h_1$ con $|h_1|\leq\alpha/3$ en $X$ y error a lo sumo $2\alpha/3$ en $V$. El residuo $r_1=g-h_1|_V$ sigue siendo continuo. Aplicamos el mismo lema al residuo con cota $\alpha_1=2\alpha/3$ y obtenemos $h_2$ con
\[
 |h_2|\leq\frac\alpha3\frac23,
 \qquad |g-h_1-h_2|\leq\alpha(2/3)^2\text{ en }V.
\]
Inductivamente, si $h_1,\ldots,h_{n-1}$ ya están construidas, aplicamos el lema al residuo continuo $g-\sum_{k=1}^{n-1}h_k|_V$, cuya cota es $\alpha(2/3)^{n-1}$. Esto produce $h_n$ continua en $X$ con
\[
 |h_n(x)|\leq\frac\alpha3(2/3)^{n-1}\quad(x\in X),
 \qquad
 \left|g(x)-\sum_{k=1}^nh_k(x)\right|\leq\alpha(2/3)^n\quad(x\in V).
\]
La suma del error es \emph{finita} en cada etapa; ése es el residuo al que se aplica el siguiente paso. Como
\[
 \sum_{n=1}^\infty\frac\alpha3(2/3)^{n-1}=\alpha,
\]
el criterio M asegura que $f=\sum_nh_n$ converge uniformemente en $X$. Sus sumas parciales son continuas y el teorema~\ref{ch12:limite-continuo} muestra que $f$ es continua. La segunda estimación, al dejar $n\to\infty$, da $f|_V=g$. La primera, sumada, da $|f|\leq\alpha$ en todo $X$.

\textbf{Valores en un intervalo.} Si $a=b$, $g$ es constante y la misma constante la extiende. Si $a<b$, aplicamos el caso anterior a
\[
 g_0=g-\frac{a+b}{2},\qquad |g_0|\leq\frac{b-a}{2}.
\]
Su extensión $f_0$ con la misma cota produce $f=f_0+(a+b)/2$, cuyos valores están en $[a,b]$ y que coincide con $g$ en $V$.

\textbf{Caso no acotado.} La función $h=\arctan g$ es continua en $V$ y toma valores estrictamente entre $-\pi/2$ y $\pi/2$. Por el caso acotado tiene una extensión continua $H:X\to[-\pi/2,\pi/2]$. No podemos aplicar directamente $\tan H$: fuera de $V$, la extensión podría alcanzar los extremos, donde la tangente no está definida. Para evitarlo, definimos
\[
 w(x)=\frac1{1+\dist(x,V)},\qquad \widetilde H(x)=w(x)H(x).
\]
Como $V$ es cerrado y no vacío, $w$ es continua, vale $1$ en $V$ y satisface $0<w(x)<1$ fuera de $V$. En $V$, $|H|<\pi/2$ porque $H=h$; fuera de $V$, el factor $w<1$ garantiza $|\widetilde H|<\pi/2$. Por ello
\[
 f(x)=\tan(\widetilde H(x))
\]
es una función real continua en todo $X$, y en $V$ cumple $f=\tan(\arctan g)=g$. Esto completa el caso que requiere especial cuidado con los extremos del intervalo.
\end{proof}

\begin{ejemplo}[Por qué se necesita un cerrado]
La función $g:(0,\infty)\to\R$, $g(x)=\cos(1/x)$, es continua y está acotada por uno. Sin embargo, para
\[
 x_n=1/(2\pi n),\qquad y_n=1/((2n+1)\pi)
\]
ambas sucesiones tienden a cero y $g(x_n)=1$, $g(y_n)=-1$. Una extensión continua a $\R$ tendría que tener en cero dos límites distintos. La condición de que el dominio original sea cerrado no puede simplemente omitirse.
\end{ejemplo}

\section{Ejercicios resueltos}
\begin{ejercicio}[Lectura 12, ejercicio 1]
Sin usar el lema de Urysohn, separe dos cerrados disjuntos $A,B\subset X$ por abiertos disjuntos que los contengan.
\end{ejercicio}
\begin{solucion}
Si alguno es vacío, elegimos para él el abierto vacío y para el otro el abierto $X$. Si ambos son no vacíos, tomamos
\[
 U=\{x:\dist(x,A)<\tfrac12\dist(x,B)\},\qquad
 W=\{x:\dist(x,B)<\tfrac12\dist(x,A)\}.
\]
Son abiertos porque son conjuntos de positividad o negatividad estricta de diferencias de funciones continuas. Si $a\in A$, entonces $\dist(a,A)=0$ y $\dist(a,B)>0$, por ser $B$ cerrado y disjunto de $A$; de ahí $a\in U$. El mismo argumento da $B\subset W$. Un punto en ambos abiertos satisfaría $\dist(x,A)<\frac14\dist(x,A)$, imposible para un número no negativo. Así $U\cap W=\varnothing$.
\end{solucion}

\begin{ejercicio}[Lectura 12, ejercicio 2]\label{ch12:soporte}
Defina el soporte de una función por
\[
 \supp(f)=\cl{\{x\in X:f(x)\ne0\}}.
\]
Sean $V$ cerrado y $U$ abierto con $V\subset U$.
\begin{enumerate}[label=(\roman*)]
\item Construya $\varphi:X\to[0,1]$ continua, igual a $1$ en $V$, y con $\supp(\varphi)\subset U$.
\item Examine la afirmación de que toda $g:V\to[a,b]$ continua tiene una extensión continua $h:X\to[a,b]$ con soporte en $U$, y resuelva la versión válida.
\end{enumerate}
\end{ejercicio}
\begin{solucion}
\textbf{(i)} Si $V=\varnothing$, tomamos $\varphi=0$; si $U=X$, tomamos $\varphi=1$. En otro caso aplicamos Urysohn a los cerrados disjuntos $V$ y $X\setminus U$ y obtenemos $u:X\to[0,1]$ con $u=0$ en $V$ y $u=1$ en $X\setminus U$. Definimos
\[
 \varphi(x)=\max\{0,1-3u(x)\}.
\]
Es continua, vale uno en $V$, y sus puntos no nulos están en $\{u<1/3\}$. Al tomar clausuras,
\[
 \supp(\varphi)\subset\{u\leq1/3\}\subset U,
\]
pues $u=1$ fuera de $U$. El margen entre $1/3$ y $1$ es lo que garantiza la inclusión para la \emph{clausura} del conjunto de no anulación, no sólo para ese conjunto.

\textbf{(ii)} El enunciado requiere una corrección cuando $0\notin[a,b]$ y $U\ne X$. Por ejemplo, si $X=\R$, $V=\{0\}$, $U=(-1,1)$ y $g(0)=1\in[1,2]$, toda función $h:\R\to[1,2]$ es no nula en todo $\R$ y su soporte es $\R$, imposible de contener en $U$. El texto original tampoco hace explícita la igualdad $h|_V=g$, que aquí incluimos para expresar el problema de extensión.

La versión válida conserva todas las propiedades si $0\in[a,b]$. Por Tietze, extendemos $g$ a $\widetilde g:X\to[a,b]$. Tomamos la función $\varphi$ del inciso (i) y ponemos $h=\varphi\widetilde g$. Es continua y coincide con $g$ en $V$. Como $0\leq\varphi\leq1$ y el intervalo contiene a cero, $h(x)$ es una combinación convexa de $\widetilde g(x)$ y cero y permanece en $[a,b]$. Además, $\{h\ne0\}\subset\{\varphi\ne0\}$ y, por clausura, $\supp(h)\subset\supp(\varphi)\subset U$.

Sin exigir $0\in[a,b]$, el mismo procedimiento da siempre una extensión con soporte en $U$ y valores en $[\min\{a,0\},\max\{b,0\}]$. Si se desea conservar exactamente el intervalo original, otra formulación válida fija un $c\in[a,b]$ y exige soporte de $h-c$ en $U$: la función $h=c+\varphi(\widetilde g-c)$ permanece en $[a,b]$, coincide con $g$ en $V$ y vale $c$ fuera del soporte de $\varphi$.
\end{solucion}

\begin{ejercicio}[Lectura 12, ejercicio 3]
Sean $U_1,\ldots,U_n$ abiertos que cubren $X$. Construya funciones continuas $\varphi_i:X\to[0,1]$ con $\supp(\varphi_i)\subset U_i$ y $\sum_{i=1}^n\varphi_i=1$.
\end{ejercicio}
\begin{solucion}
Si $X=\varnothing$, las únicas funciones con ese dominio satisfacen todas las condiciones. Supongamos $X\ne\varnothing$ y $n\geq1$. Para cada $i$ definimos
\[
 r_i(x)=\begin{cases}
 \min\{1,\dist(x,X\setminus U_i)\},&U_i\ne X,\\
 1,&U_i=X.
 \end{cases}
\]
Son continuas y no negativas. Si $U_i\ne X$, el complemento es cerrado no vacío y $r_i(x)>0$ exactamente cuando $x\in U_i$; esta afirmación también vale para $U_i=X$. Por la propiedad de cubierta, $S(x)=\sum_i r_i(x)>0$ en todo punto. Pongamos $w_i=r_i/S$. Son continuas, no negativas y suman uno.

No basta tomar directamente $w_i$ como partición subordinada en el sentido del soporte: su conjunto no nulo está en $U_i$, pero su clausura podría alcanzar la frontera. Introducimos un margen:
\[
 V_i=\{x:w_i(x)>1/(2n)\},\qquad
 g_i(x)=\max\{w_i(x)-1/(2n),0\}.
\]
Para cada $x$, algún $w_i(x)$ es al menos $1/n$, ya que la suma es uno. Por ello los $V_i$ cubren $X$ y $G=\sum_i g_i$ es estrictamente positiva. Además,
\[
 \cl{V_i}\subset\{w_i\geq1/(2n)\}\subset\{w_i>0\}=U_i.
\]
Ésta realiza explícitamente la reducción de la cubierta sugerida en la lectura original. Finalmente, definimos $\varphi_i=g_i/G$. Son continuas, están entre cero y uno y suman uno. Su conjunto de no anulación es $V_i$, luego $\supp(\varphi_i)=\cl{V_i}\subset U_i$. Hemos construido una partición finita de la unidad con todas las inclusiones de soporte justificadas.
\end{solucion}
\index{particion de la unidad@partición de la unidad}\index{soporte de una funcion@soporte de una función}

\section{Extensión, localización y suma}
Las funciones de separación permiten fijar valores sobre cerrados; las series uniformemente convergentes convierten aproximaciones en una extensión exacta; y las funciones de soporte controlado permiten localizar esa extensión. Una partición de la unidad combina estas ideas para repartir una construcción global entre distintos abiertos sin perder continuidad.
